\chapter[Discussion]{Discussion}
\label{Chap:Discussion}

% \newcommand{\mhz}{\>\mathrm{MHz}}
% \newcommand{\ghz}{\>\mathrm{GHz}}

\section{Interpretation of results}
The interpretation of RXJ1720+2638 has previously been uncertain. An 8.44 GHz measurement suggested tentative evidence for steepening relative to lower frequencies \citep{Giacintucci2014a}, whereas \textit{uv}-plane modelling at 15.5 GHz was consistent with a single power-law spectrum \citep{Perrott2021}. The results of this thesis provide direct observational evidence in support of a single power-law spectrum up to at least 10 GHz. This is consistent with the conclusions of \citet{Perrott2021}, who argued that incomplete \textit{uv}-coverage at 8.44 GHz artificially reduced the measured flux density. 

The flux density measured over the tail region of the RXJ1720+2638 minihalo is approximately three times higher than that predicted by extrapolation of previous measurements. One possible explanation is the use of a different spatial mask, which may include a larger fraction of the brighter central emission within the region classified as the tail. As the mask was estimated by eye from the previous studies, this is plausible. The detectable high-frequency emission extends over only approximately half of the masked tail region. This reduction in spatial extent with increasing frequency is consistent with a steepening spectral index along the tail, as radiative losses progressively suppress the population of high-energy electrons. The results for this cluster will be explored in more depth in a forthcoming study (Colquhoun et al., in prep). 

The findings for RXCJ1115+0129 are somewhat less constraining, primarily due to the limited number of flux density measurements for this minihalo. The previous study by \citet{Trehaeven2023} found a spectral index of $\alpha=1.0\pm1.1$ between 1--1.5 GHz, with the uncertainty dominated by subtraction errors. The results presented in this thesis are consistent with their mean value and reduce the uncertainty by approximately a factor of six, although the uncertainty remains dominated by source subtraction. The northern and southern spatial steepening reported in the previous study was not explored in depth in this thesis. However, the narrower observed spatial extent at 10 GHz is near-aligned with the direction of their measured steepening. A spatially resolved point-to-point analysis combining the 1.3 GHz MeerKAT data with the 10 GHz observations presented here would provide a more robust test for this feature.

Several previous flux density measurements for the RXJ2129+0005 minihalo were available for comparison with the value measured in this thesis. The 10 GHz flux density is found to be slightly higher than expected from the extrapolation of the previous studies, although the $1\sigma$ uncertainties overlap. The excess flux may arise from two regions of emission near the masking boundary, including the unsubtracted source S2 and a diffuse region to the east (see Figure \ref{fig:rxj2129_sub_image}). These may contaminate the masked area and introduce flux not associated with the minihalo. The complex masking region defined by \citet{Giacintucci2014} was approximated here by an ellipse, which further modifies the measured flux density. Notably, the flux density measured within the $3\sigma$ contour is less than half that of the larger mask ($S_{3\sigma} = 0.15\pm0.1$ mJy compared to $S_{mask} = 0.41\pm0.1$ mJy), which demonstrates the sensitivity of the integrated flux to the adopted mask.

The final cluster measured in this thesis, A478, presents a more intriguing case. Previous studies placed only an upper limit on the spectral index of $\alpha<1$, whereas the flux density measured in this thesis implies a far greater spectral index of $\alpha = 2.02\pm0.10$. The in-band spectral index is even steeper at $\alpha = 3.34\pm0.18$, suggesting significant spectral curvature between 1.4 and 10 GHz. This result should be interpreted with caution, however, as it is based upon only two direct measurements and one non-detection. The adopted mask and \textit{uv}-range likely differ from those used by \citet{Giacintucci2014} and \citet{Savini2019}. Spatial filtering at 10 GHz may result in missing flux, analagous to the 8.44 GHz measurement of RXJ1720+2638. Additional intermediate-frequency observations (e.g. 144 MHz to 5 GHz) are required to determine whether this apparent steepening reflects genuine particle ageing or residual systematic effects. Nevertheless, A478 remains an outlier compared to the other clusters in this sample.

The clusters RXJ1720+2638, RXCJ1115+0129, and RXJ2129+0005 exhibit total 10 GHz flux densities consistent, within uncertainties, with the extrapolated power-law spectra derived from lower-frequency measurements. The absence of clear high-frequency curvature in the integrated spectra suggests that strong radiative ageing effects are not dominant over this frequency range. This behaviour is broadly consistent with the expectations from the hadronic scenario, in which secondary electrons are continuously injected and maintain a single power-law energy distribution to high frequencies.

In contrast, the turbulent acceleration models predict spectral steepening at high-frequencies once the stochastic acceleration becomes inefficient at compensating for radiative losses of the CRe population. The results presented here therefore suggest that turbulent reacceleration, if present, is not dominant at the minihalo scale or that an additional mechanism is sustaining the high-energy CRe population. However, the spatially resolved steepening observed in the tail of RXJ1720+2638, together with the tentative evidence of curvature in A478, indicates that multiple mechanisms may operate within individual clusters. This interpretation is consistent with previous studies (e.g \citealp{Perrott2021, Riseley2024}).

\section{Incomplete clusters}
For the majority of the remaining clusters in this sample a minihalo is detected at 10 GHz. However, difficulties encountered during the subtraction of contaminating emission prevented robust estimation of the minihalo flux in several cases. These difficulties can be categorised as arising from (i) variability of compact sources between observing epochs and (ii) residual diffuse emission not associated with the minihalo in the subtracted image. 

Variability of compact sources between observations at different epochs presents a significant challenge for the self-calibration and source subtraction. Point source modelling was performed independently for each dataset in order to mitigate this effect. However, this did not fully resolve the residual artefacts in the individual models nor the joint image. The C-configuration observations were typically the dominant contributor to subtraction errors in the combined datasets image, although the precise reason behind this remains unclear. Possible factors are the \texttt{tclean} parameters for imaging of the point sources and residual calibration errors.

The clusters where a minihalo was not detected at 10 GHz are A1795 and A2626. In both cases, the absence of a clear detection is attributed to imaging artefacts obscuring potential diffuse emission, rather than insufficient surface brightness sensitivity. These clusters contain bright central AGNs that complicated the self-calibration process. It is plausible that with additional refinement of the self-calibration process, a reliable image of the minihalo emission could be recovered. 

The models used for source subtraction were constructed using a fixed \textit{uv}-cut and imaging parameters derived from an analysis of a single cluster, RXJ1720+2638 (see Section \ref{sec:source_subtraction}). Although manual inspection of each resulting image was performed, these parameters likely require optimisation on a per-cluster basis. A key systematic limitation arises because the adopted \textit{uv}-cut corresponds to a constant angular scale on the sky, rather than a constant physical scale.

The clusters in this sample span physical scales between 0.72 and 7.5 kpc arcsec$^{-1}$, increasing with redshift under the adopted cosmology. Consequently, the same angular \textit{uv}-cut corresponds to substantially different physical scales across the sample. RXJ1720+2638 lies at an intermediate redshift within the cluster sample. For lower-redshift clusters an identical \textit{uv}-cut selects smaller physical scales. In such cases, diffuse emission from AGN lobes may not be fully included in the subtraction model, leading to residual contamination in the subtracted image (e.g. 2A0335+096). 

Conversely, for higher-redshift clusters, the same \textit{uv}-cut corresponds to larger physical scales. This increase the risk that extended minihalo emission is partially incorporated into the subtraction model, resulting in over-estimation of compact source flux densities, and subsequently, over-subtraction in the final image (e.g ACTJ0022-0036). 

The impact of adopting a constant angular \textit{uv}-cut is dependent on the minihalo extent, cluster redshift, and the morphology of embedded radio sources. A physically motivated \textit{uv}-cut scaled to the observed extent may reduce this systematic bias in future analyses.

\section{Measurement Limitations}
There are several fundamental observational limitations affecting the detectability of minihalos in this study. The primary constraints arise from the maximum recoverable scale, surface brightness sensitivity, and angular resolution of the interferometer configuration. 

\subsection{Maximum recoverable scale}
The largest angular scale that the interferometer can recover is set by the minimum projected baseline length. For these VLA observations, the shortest baselines correspond to 1 k$\lambda$. Approximating the angular scale as $\theta \approx 1/\vec u$ yields $10^{-3}$ rad or $\sim\!3.44$ arcminutes. 

For a characteristic minihalo diameter of $\sim\!300$ kpc, this angular scale corresponds to a minihalo located at redshift $z \sim 0.07$. Below this redshift, emission at scales comparable to or larger than this minihalo diameter may be partially resolved out, leading to an underestimation of the total flux density. 

This effect is most relevant for three clusters in this sample, 2A0335+096, A2626, and A1795. Of these clusters, only 2A0335+096 has a source-subtracted image. For this cluster $3.44'$ corresponds to $\approx150$ kpc. Diffuse emission extending beyond this scale will not be fully recovered. However, the lowest-redshift cluster for which a total 10 GHz flux density was successfully measured is RXJ1720+2638 at $z=2.775$. This corresponds to a maximum physical scale of $\sim\!570$ kpc, suggesting that spatial filtering is unlikely to significantly bias the integrated flux densities reported here. Matching the \textit{uv}-coverage with prior studies, as performed for RXJ1720+2638, also ensures any spatial filtering is consistent between studies.

\subsection{Surface brightness sensitivity}

The surface brightness sensitivity is limited by image noise, which arises from residual RFI, calibration uncertainties, and instrumental thermal noise. For the 10 GHz observations presented here, the typical rms noise level is approximately $3$ $\mu$Jy beam$^{-1}$. 

For steep-spectrum or intrinsically low-brightness minihalos, emission may fall below the adopted $3\sigma$ detection threshold. This can artifically reduce the apparent spatial extent of the emission relative to lower-frequency observations, where surface brightness is higher. In such cases, only lower limits on the spectral index can be derived. As discussed previously, spatial masks defined in earlier studies were adopted where available to ensure consistency in flux density measurements. 

One method to increase the signal-to-noise ratio of diffuse emission is to smooth an image, typically by convolution with a Gaussian kernel. Such spatial averaging enhances low-surface brightness emission but reduces spatial resolution. This process disproportionally affects compact sources relative to extended emission, as their flux is redistributed over a larger area. Consequently, smoothing is more appropriately applied to source-subtracted images. This approach was not explored in this thesis and remains a potential avenue for future refinement of these measurements.

\subsection{Angular resolution}
The angular resolution, defined by the restoring beam, has the greatest impact on the accuracy of the source subtraction. For the 10 GHz observations, the typical restoring beam is approximately $4"\times3"$. Higher angular resolution improves the modelling of compact sources and their associated lobes or jets, enabling more effective separation from the diffuse minihalo emission.

Many compact sources are unresolved at this beam size. Improved resolution allows subtraction to occur over smaller areas, thereby reducing subtraction uncertainties. Since the subtraction error scales with the area over which the model is removed, $\sigma_{sub} \propto A_{sub}$, improved resolution directly reduces residual contamination. As discussed previously, the inclusion of C-configuration data provides the necessary angular resolution in the joint imaging process.

\section{SZ effect}
The SZ effect represents a potential contaminant to the flux densities measured in this thesis. As it arises from IC-scattering through the cluster, it is naturally a cluster-scale phenomenon extending well beyond the spatial extent of the minihalo. Such a large-scale structure may be partially sampled at the shortest baselines of the 10 GHz observations. 

The thermal SZ spectrum has a null frequency at 217 GHz, where there is no contribution. Below this frequency the effect is negative, suppressing the measured flux density, while above it the effect is positive. At 10 GHz the effect is therefore negative. However, the amplitude of the SZ decrement is expected to be small. Figure \ref{fig:sz_amp} depicts example spectra for the SZ effect.

\begin{figure}
    \includegraphics[width=0.85\textwidth]{Chapter6/sz_ampl.pdf}
    \caption{Thermal SZ intensity spectrum illustrating characteristic frequency dependence. Curves correspond to representative Compton-$y$ values, with larger $y$ typically associated with more massive clusters. The dotted black line denotes the null frequency at 217 GHz.}
    \label{fig:sz_amp}
\end{figure}


\citet{Perrott2021} estimated the maximum SZ signal for RXJ1720+2638 for the VLA observation at 8.44 GHz. They found a negative maximum amplitude of approximately 80 $\mu$Jy at the shortest baselines. Given the similar observing frequency and comparable interferometer configuration, this provides a reasonable estimate for the expected SZ signal. For RXJ1720+2638 this is negligible compared to the measured minihalo flux and lies well within the quoted uncertainty.

The remaining clusters where minihalo fluxes are obtained (A478, RXCJ1115+0129, and RXJ2129+0005) have masses within $-26\%$ to $+18\%$ of RXJ1720+2638. As the SZ effect scales primarily with the integrated thermal pressure (and therefore cluster mass), a very conservative upper bound of $<0.1$ mJy is considered for all clusters. This is within the measurement uncertainties for RXCJ1115+0129 and RXJ2129+0005. For A478, this represents approximately 33\% of the measured flux. However, this difference is insufficient to reconcile the order of magnitude discrepancy to the lower limit on the spectral index reported by \citet{Savini2019}.

Given that the SZ contribution is comparable to or below the measurement uncertainties, it is unlikely to significantly affect the conclusions presented here. A more rigorous treatment would involve modelling the SZ signal for each cluster using gas pressure profiles and propagating through the interferometer response, as performed by \citet{Perrott2021}. Such analysis is beyond the scope of this thesis, but would enable the direct correction of the minihalo flux densities.
